\documentclass[../../../include/open-logic-section]{subfiles}

\begin{document}

\olfileid{sth}{ordinals}{basic}
\olsection{الخصائص الأساسية للأعداد الترتيبية}

قد تبين أن أوائل الأعداد الترتيبية هي الأعداد الطبيعية. وإنما نبني
نظرية الأعداد الترتيبية، في المقام الأول، لنتجاوز بمبدأ الاستقراء
مجاله في الأعداد الطبيعية. ونمهد لذلك بنتائج أولية متتابعة.

\begin{lem}\ollabel{ordmemberord}
كل !!{element} من عدد ترتيبي عدد ترتيبي أيضًا.
\end{lem}

\begin{proof}
نأخذ عددًا ترتيبيًا~$\alpha$، ونفرض~$b \in \alpha$. فتعدي~$\alpha$
يعطي~$b \subseteq \alpha$؛ ولذلك ترتب~$\in$ المجموعة~$b$ ترتيبًا
جيدًا، إذ إن~$\in$ ترتب~$\alpha$ ترتيبًا جيدًا.

بقي تعدي~$b$. لنفرض~$x \in c \in b$؛ فيلزم~$c \in \alpha$ من
$b \subseteq \alpha$. وبتعدي~$\alpha$ مرة أخرى نحصل على~$c \subseteq \alpha$،
ثم~$x \in \alpha$. وهكذا~$x, c, b \in \alpha$. ولما كانت~$\in$
ترتب~$\alpha$ ترتيبًا جيدًا، كانت~$\in$ علاقة متعدية على~$\alpha$،
بموجب~\olref[sth][ordinals][wo]{wo:strictorder}. فمن~$x \in c \in b$
يلزم~$x \in b$؛ وبالتعميم~$c \subseteq b$.
\end{proof}

\begin{cor}\ollabel{ordissetofsmallerord}
لدينا~$\alpha = \Setabs{\beta \in \alpha}{\beta \text{ عدد ترتيبي}}$
لكل عدد ترتيبي~$\alpha$.
\end{cor}

\begin{proof}
هذه نتيجة مباشرة لـ\olref{ordmemberord}.
\end{proof}

وسنثبت نتيجتين أساسيتين، هما~\olref{ordinductionschema}
و\olref{ordtrichotomy}؛ ومعناهما الإجمالي أن الانتماء يرتب الأعداد
الترتيبية نفسها ترتيبًا جيدًا.

\begin{thm}[الاستقراء المتناهي التجاوز]\ollabel{ordinductionschema}
لكل صيغة~$\phi(x)$ يصح:
\[
\text{إذا كان }\exists \alpha \phi(\alpha)\text{، فإن }
\exists \alpha(\phi(\alpha)
\land (\forall \beta \in \alpha) \lnot \phi(\beta))
\]
والمكممات الظاهرة هنا مقيدة بالأعداد الترتيبية وإن لم يُصرَّح بالقيد.
\end{thm}

\begin{proof}
%\olref{ordinductionschema1}.
لنفرض~$\phi(\alpha)$ لعدد ترتيبي~$\alpha$. فإن كان
$(\forall \beta \in \alpha) \lnot \phi(\beta)$ تم المطلوب. وإلا وضعنا
$X = \Setabs{\beta \in \alpha}{\phi(\beta)}$، فكان~$X \neq \emptyset$.
وبما أن~$\in$ ترتب~$\alpha$ ترتيبًا جيدًا، فلـ$X$ !!{element} أصغر
بحسب~$\in$، نرمز إليه بـ$\gamma$. ومن تعريف~$X$ نجد~$\phi(\gamma)$؛
ثم إن~$\gamma$ عدد ترتيبي بمقتضى~\olref{ordmemberord}.
فإذا كان~$\beta \in \gamma$ لزم~$\beta \in \alpha$ من تعدي~$\alpha$،
وامتنع~$\beta \in X$ لأصغرية~$\gamma$. إذن~$\lnot\phi(\beta)$،
وبالتعميم~$(\forall \beta \in \gamma)\lnot\phi(\beta)$.
%\olref{ordinductionschema2}. Suppose there is some ordinal
% $\gamma$ such that $\lnot\phi(\gamma)$. Then by
% \olref{ordinductionschema1}, there is an $\in$-minimal ordinal
% $\alpha$ for which $\lnot\phi(\alpha)$. So $(\forall \beta <
% \alpha) \phi(\beta)$, rendering the antecedent of the conditional
% false.
\end{proof}
\noindent
ولـ\olref{ordinductionschema} صياغة مكافئة بالمخطط الآتي:
\[
\text{إذا كان }\forall \alpha((\forall \beta \in \alpha)\phi(\beta) \lif
\phi(\alpha))\text{، فإن }\forall \alpha\phi(\alpha)
\]
فما علينا إلا أن نعوض بـ$\lnot\phi(\alpha)$ في~\olref{ordinductionschema}،
ثم نجري التحويلات المنطقية الأولية.

\begin{thm}[التثليث]\ollabel{ordtrichotomy}
لدينا~$\alpha \in \beta \lor \alpha = \beta \lor \beta \in \alpha$
لكل عددين ترتيبيين~$\alpha$ و$\beta$.
\end{thm}

\begin{proof}
نستعمل الاستقراء مرتين، أي نطبق~\olref{ordinductionschema} تطبيقين.
ونقول إن~$x$ \emph{قابل للمقارنة} مع~$y$ إذا وفقط إذا كان
$x \in y \lor x = y \lor y \in x$.

نفرض أولًا أن كل عدد ترتيبي في~$\alpha$ قابل للمقارنة مع \emph{كل}
عدد ترتيبي. ثم نفرض، لاستقراء ثان، أن~$\alpha$ قابل للمقارنة مع كل
عدد ترتيبي في~$\beta$. والمطلوب الآن مقارنة~$\alpha$ بـ$\beta$.
فإذا ثبتت، دل الاستقراء على~$\beta$ على أن~$\alpha$ قابل للمقارنة
مع كل عدد ترتيبي، ثم دل الاستقراء على~$\alpha$ على أن \emph{كل}
عدد ترتيبي قابل للمقارنة مع \emph{كل} عدد ترتيبي. ويكفينا لهذا أن
نفرض~$\alpha \notin \beta$ و$\beta \notin \alpha$ ونستنتج~$\alpha = \beta$.

نثبت أولًا~$\alpha \subseteq \beta$. خذ~$\gamma \in \alpha$؛ فهو
عدد ترتيبي بمقتضى~\olref{ordmemberord}، ولذلك يقبل~$\gamma$ المقارنة
مع~$\beta$ بفرض الاستقراء الأول. فإن كان~$\gamma = \beta$ أو
$\beta \in \gamma$، لزم~$\beta \in \alpha$، مع الاستعانة بتعدي~$\alpha$
حيث يلزم. وهذا خلاف الفرض، فلا يبقى إلا~$\gamma \in \beta$.
وبالتعميم نحصل على~$\alpha \subseteq \beta$.

وبمثل هذه الحجة تمامًا، ولكن بفرض الاستقراء الثاني، يثبت
$\beta \subseteq \alpha$؛ فينتج~$\alpha = \beta$.
\end{proof}\noindent
لهذا نكتب أحيانًا~$\alpha <\beta$ مكان~$\alpha \in \beta$، فقد
قامت~$\in$ مقام علاقة ترتيب. وليس وراء هذا الاصطلاح أمر عميق:
إنما هو مألوف، وكتابة~$\alpha \leq \beta$ أيسر من كتابة
$\alpha \in \beta \lor \alpha = \beta$.
\footnote{يجوز من حيث الإمكان أن نكتب~$\alpha \mathrel{\underline{\in}} \beta$،
إلا أنه ترميز خارج عن الاصطلاح المعتاد تمامًا.}

ويتفرع عما تقدم نتيجتان قريبتان، يظهر في أولاهما استعمال الترميز الجديد.

\begin{cor}\ollabel{ordordered}
إذا كان~$\exists \alpha\phi(\alpha)$، لزم
$\exists \alpha(\phi(\alpha) \land \forall \beta(\phi(\beta) \lif \alpha \leq \beta))$.
وكذلك، لكل أعداد ترتيبية~$\alpha, \beta, \gamma$، يصح
$\alpha \notin \alpha$ ويصح~$\alpha \in \beta \in \gamma \lif \alpha \in \gamma$.
\end{cor}

\begin{proof}
يجري البرهان على مثال~\olref[wo]{wo:strictorder} بعينه.
\end{proof}

\begin{prob}
استوفِ الحجة التي قيل عنها «بمثل هذه الحجة تمامًا» في برهان
\olref[sth][ordinals][basic]{ordtrichotomy}.
\end{prob}

\begin{cor}\ollabel{corordtransitiveord}
تكون~$A$ عددًا ترتيبيًا إذا وفقط إذا كانت~$A$ مجموعة متعدية عناصرها
أعداد ترتيبية.
\end{cor}

\begin{proof}
أما الاتجاه \emph{من اليسار إلى اليمين} فمن~\olref{ordmemberord}.
وأما \emph{من اليمين إلى اليسار}، فإذا كانت~$A$ مجموعة متعدية من
الأعداد الترتيبية، فإن~$\in$ ترتب~$A$ ترتيبًا جيدًا بمقتضى
\olref{ordinductionschema} و\olref{ordtrichotomy}.
\end{proof}

قد فسرنا~\olref{ordinductionschema} و\olref{ordtrichotomy} بأن
$\in$ ترتب الأعداد الترتيبية ترتيبًا جيدًا. لكن هذا القول يحتاج إلى
\emph{احتراز شديد}، كما يتبين من النتيجة التالية.

\begin{thm}[مفارقة بورالي--فورتي]\ollabel{buraliforti}
يمتنع أن تكون هناك مجموعة تضم جميع الأعداد الترتيبية.
\end{thm}

\begin{proof}
لنفترض خلاف ذلك، ولتكن~$O$ مجموعة الأعداد الترتيبية كلها. إذا كان
$\alpha \in \beta \in O$، كان~$\alpha$ عددًا ترتيبيًا بحسب
\olref{ordmemberord}، فيلزم~$\alpha \in O$. وعليه تكون~$O$ متعدية،
فتكون عددًا ترتيبيًا بمقتضى~\olref{corordtransitiveord}؛ وينتج
$O \in O$، خلافًا لـ\olref{ordordered}.
\end{proof}
\noindent
نُسبت النتيجة إلى~\citeauthor{Burali-Forti1897}. غير أن كانتور، في
رسالته إلى ددكند سنة~1899، كان أول من أدرك بوضوح \emph{التناقض}
الناشئ من افتراض مجموعة لجميع الأعداد الترتيبية. ويبين فان هاينورت ذلك بقوله:
\begin{quote}
أما بورالي--فورتي نفسه فعدّ التناقض برهانًا، بطريق \emph{الخلف}،
على أن الترتيب الطبيعي للأعداد الترتيبية ليس إلا ترتيبًا جزئيًا.
\citep[p.~105]{Heijenoort1967}
\end{quote}
وبصرف النظر عن خطئه، حاصل ما تقدم أن الأعداد الترتيبية مجموعات
يرتب الانتماء كل واحدة منها ترتيبًا جيدًا، ويرتبها جميعًا ترتيبًا
جيدًا أيضًا، من غير أن يكون مجموعها مجموعة.

ونختم بخواص أولية أخرى تتفرع عن~\olref{ordinductionschema}
و\olref{ordtrichotomy}.

\begin{prop}
لا تكون المتتالية المتناقصة تمامًا من الأعداد الترتيبية إلا متناهية.
\end{prop}

\begin{proof}
لو وجدت متتالية غير متناهية متناقصة تمامًا
$\alpha_0 > \alpha_1 > \alpha_2 > \ldots$، لما كان لها عنصر أصغري
بحسب~$<$؛ وذلك يناقض~\olref{ordinductionschema}.
\end{proof}

\begin{prop}\ollabel{ordinalsaresubsets}
لكل عددين ترتيبيين يصح~$\alpha \subseteq \beta \lor \beta \subseteq \alpha$،
حيث~$\alpha, \beta$ هما العددان.
\end{prop}

\begin{proof}
إن كان~$\alpha \in \beta$، لزم~$\alpha \subseteq \beta$ لتعدي~$\beta$.
وإن كان~$\beta \in \alpha$، لزم كذلك~$\beta \subseteq \alpha$.
أما إن كان~$\alpha = \beta$، فيصح~$\alpha \subseteq \beta$
و$\beta \subseteq \alpha$ معًا. وقد حصرت~\olref{ordtrichotomy}
جميع الحالات في هذه الثلاث، فتم المطلوب.
\end{proof}

\begin{prop}\ollabel{ordisoidentity}
لدينا~$\alpha = \beta$ إذا وفقط إذا كان~$\ordeq{\alpha}{\beta}$،
وذلك لكل عددين ترتيبيين~$\alpha, \beta$.
\end{prop}

\begin{proof}
الأعداد الترتيبية ترتيبات جيدة؛ فتنتج الدعوى من التثليث
(\olref[basic]{ordtrichotomy}) ومن~\olref[iso]{wellordnotinitial} مباشرة.
\end{proof}

\begin{prob}\ollabel{probunionordinalsordinal}
برهن أنه متى كانت عناصر~$X$ كلها أعدادًا ترتيبية، كان~$\bigcup X$
عددًا ترتيبيًا أيضًا.
\end{prob}

\end{document}
